\chapter{Type System and Type Checking}
\label{chap:types}

\section{Bidirectional Type Inference Architecture}
\label{sec:types:bidirectional}

NumLang combines the expressive simplicity of Hindley-Milner parametric polymorphism with the deterministic predictability of bidirectional type checking \citep{pierce2000local}. Implemented in \texttt{src/typecheck/checker.rs}, the typechecker enforces strict static type safety prior to intermediate representation lowering.

The type system is organized around two complementary operational judgments:
\begin{enumerate}
    \item \textbf{Type Synthesis} ($\Gamma \vdash e \Rightarrow \tau$): Given an expression $e$ under typing context $\Gamma$, the type checker infers its unique principal type $\tau$.
    \item \textbf{Type Checking} ($\Gamma \vdash e \Leftarrow \tau$): Given an expression $e$ and an expected type $\tau$, the type checker verifies that $e$ conforms to $\tau$, propagating type information downward into sub-expressions.
\end{enumerate}

The typing environment $\Gamma$ maps term variables to type schemes $\sigma = \forall \vec{\alpha}.\,\tau$. The core primitive types comprise 64-bit signed and unsigned integers (\texttt{i64}, \texttt{u64}), sub-word integers (\texttt{i32}, \texttt{i16}, \texttt{i8}, \texttt{u32}, \texttt{u16}, \texttt{u8}), IEEE-754 floating-point numbers (\texttt{f64}, \texttt{f32}), booleans (\texttt{bool}), unit (\texttt{()}), heap-allocated boxes (\texttt{Box<T>}), fixed-size arrays (\texttt{[T; N]}), algebraic data types (enums), user-defined record structures (structs), and first-class function signatures (\texttt{fn(T1, ...) -> T2}).

\section{Formal Typing Rules}
\label{sec:types:rules}

We formalize the core typing rules of NumLang using standard natural deduction inference notation:

\begin{align}
\text{\bf (Var)} \quad & \frac{x : \sigma \in \Gamma \quad \tau = \text{instantiate}(\sigma)}{\Gamma \vdash x \Rightarrow \tau} \\[0.8em]
\text{\bf (Lit-Int)} \quad & \frac{}{\Gamma \vdash c_{\text{int}} \Rightarrow \mathtt{i64}} \qquad
\text{\bf (Lit-Bool)} \quad \frac{}{\Gamma \vdash c_{\text{bool}} \Rightarrow \mathtt{bool}} \\[0.8em]
\text{\bf (BinOp-Arith)} \quad & \frac{\Gamma \vdash e_1 \Leftarrow \mathtt{i64} \quad \Gamma \vdash e_2 \Leftarrow \mathtt{i64}}{\Gamma \vdash e_1 + e_2 \Rightarrow \mathtt{i64}} \\[0.8em]
\text{\bf (BinOp-Cmp)} \quad & \frac{\Gamma \vdash e_1 \Rightarrow \tau \quad \Gamma \vdash e_2 \Leftarrow \tau}{\Gamma \vdash e_1 == e_2 \Rightarrow \mathtt{bool}} \\[0.8em]
\text{\bf (Let-Val)} \quad & \frac{\Gamma \vdash e_1 \Rightarrow \tau_1 \quad \Gamma, x : \text{generalize}(\Gamma, \tau_1) \vdash e_2 \Rightarrow \tau_2}{\Gamma \vdash \mathtt{let}\; x = e_1;\; e_2 \Rightarrow \tau_2} \\[0.8em]
\text{\bf (If-Then-Else)} \quad & \frac{\Gamma \vdash e_{\text{cond}} \Leftarrow \mathtt{bool} \quad \Gamma \vdash e_1 \Rightarrow \tau \quad \Gamma \vdash e_2 \Leftarrow \tau}{\Gamma \vdash \mathtt{if}\; e_{\text{cond}}\; \{ e_1 \}\; \mathtt{else}\; \{ e_2 \} \Rightarrow \tau} \\[0.8em]
\text{\bf (Box-Alloc)} \quad & \frac{\Gamma \vdash e \Rightarrow \tau}{\Gamma \vdash \mathtt{box}(e) \Rightarrow \mathtt{Box}\langle \tau \rangle} \\[0.8em]
\text{\bf (Deref)} \quad & \frac{\Gamma \vdash e \Rightarrow \mathtt{Box}\langle \tau \rangle}{\Gamma \vdash \mathtt{deref}(e) \Rightarrow \tau} \\[0.8em]
\text{\bf (Fn-App)} \quad & \frac{\Gamma \vdash e_{\text{fn}} \Rightarrow \mathtt{fn}(\tau_1, \dots, \tau_n) \to \tau_{\text{ret}} \quad \Gamma \vdash e_i \Leftarrow \tau_i}{\Gamma \vdash e_{\text{fn}}(e_1, \dots, e_n) \Rightarrow \tau_{\text{ret}}}
\end{align}

\section{Parametric Polymorphism and Monomorphization}
\label{sec:types:monomorphize}

NumLang supports explicit generic functions parameterized over type variables:
\begin{lstlisting}[language=NumLang]
fn identity<T>(x: T) -> T {
    return x;
}
\end{lstlisting}
Rather than employing boxed runtime representations or dictionary passing, NumLang compiles generic abstractions via static \emph{monomorphization} (\texttt{src/opt/monomorphize.rs}). 

Prior to MIR lowering, the monomorphization pass analyzes the call graph originating from \texttt{main}. For each generic function invocation $f\langle \tau_1, \dots, \tau_n \rangle$, the compiler instantiates a concrete specialized clone with specialized identifier $f\_\tau_1\_\dots\_\tau_n$. Because NumLang prohibits polymorphic recursion (enforced by a cycle-detection check on polymorphic call chains), monomorphization terminates unconditionally and yields a strictly monomorphic program.

\section{Algebraic Data Types and Recursive Well-Formedness}
\label{sec:types:adts}

User-defined algebraic data types are declared via the \texttt{enum} construct:
\begin{lstlisting}[language=NumLang]
enum List {
    Nil,
    Cons(i64, Box<List>),
}
\end{lstlisting}

\subsection{Well-Formedness Criteria}
To ensure sound memory layout and prevent infinite-size value structures, the type checker enforces three structural well-formedness invariants:
\begin{enumerate}
    \item \textbf{Inductive Base Case}: Every recursive enum must contain at least one non-recursive variant (e.g., \texttt{Nil}). Enums where all variants recursively contain the enum itself are rejected as unconstructible.
    \item \textbf{Indirection via Box}: Direct recursive containment (\texttt{Cons(i64, List)}) is rejected at compile time. Recursive cycles in type definitions must be broken by an explicit indirection (\texttt{Box<T>}), guaranteeing finite pointer sizing.
    \item \textbf{Unique Tag Discriminants}: Each constructor within an enum is assigned an integer discriminant tag $0 \le \delta < k$, uniquely identifying the variant during pattern matching and MIR lowering.
\end{enumerate}

\section{Pattern Matching Exhaustiveness Checking}
\label{sec:types:exhaustiveness}

Pattern matching statements and expressions are validated for both \emph{exhaustiveness} (every possible value of the scrutinized type is matched by at least one arm) and \emph{redundancy} (no arm is shadowed by prior patterns).

NumLang implements the classic Maranget matrix-based pattern exhaustiveness algorithm \citep{maranget2007warnings}. Patterns are organized into a pattern matrix $P$, where rows correspond to match arms and columns correspond to sub-terms of the scrutinized tuple. By systematically expanding constructor signatures, the algorithm determines if the vector of wildcard patterns $(\_, \dots, \_)$ is useful relative to $P$. If useful, the compiler constructs a concrete witness value representing the unmatched case and reports it in a structured diagnostic.

\section{Diagnostic Subsystem and Error Reporting}
\label{sec:types:diagnostics}

All errors detected during type inference and exhaustiveness checking are structured into domain diagnostics (\texttt{src/diagnostic.rs}). Consider the following invalid program:
\begin{lstlisting}[language=NumLang]
fn type_error_example() -> i64 {
    let x: i64 = 42;
    let y: bool = true;
    return x + y;
}
\end{lstlisting}
The compiler emits a colorized diagnostic citing the exact source span, the inferred type (\texttt{bool}), and the expected type (\texttt{i64}):
\begin{lstlisting}
error[E0308]: mismatched types
  --> src/test.nl:4:16
   |
 4 |     return x + y;
   |                ^ expected `i64`, found `bool`
   |
   = note: binary operator `+` requires both operands to have identical numeric type
\end{lstlisting}
Zero unhandled panics occur anywhere within the type checking pipeline.

\section{Complete Natural Deduction Typing Judgments}
\label{sec:types:complete_judgments}

To establish the mathematical soundness of NumLang's bidirectional type system, we present the complete set of natural deduction inference rules for statement execution and expression typing:

\begin{align}
\text{\bf (T-While)} \quad & \frac{\Gamma \vdash e_{\text{cond}} \Leftarrow \mathtt{bool} \quad \Gamma \vdash s_{\text{body}} \Leftarrow \mathtt{()}}{\Gamma \vdash \mathtt{while}\; e_{\text{cond}}\; \{ s_{\text{body}} \} \Rightarrow \mathtt{()}} \\[0.8em]
\text{\bf (T-For-Range)} \quad & \frac{\Gamma \vdash e_{\text{start}} \Leftarrow \mathtt{i64} \quad \Gamma \vdash e_{\text{end}} \Leftarrow \mathtt{i64} \quad \Gamma, i : \mathtt{i64} \vdash s_{\text{body}} \Leftarrow \mathtt{()}}{\Gamma \vdash \mathtt{for}\; i \;\mathtt{in}\; e_{\text{start}}..e_{\text{end}}\; \{ s_{\text{body}} \} \Rightarrow \mathtt{()}} \\[0.8em]
\text{\bf (T-Array-Index)} \quad & \frac{\Gamma \vdash e_{\text{arr}} \Rightarrow [T; N] \quad \Gamma \vdash e_{\text{idx}} \Leftarrow \mathtt{i64}}{\Gamma \vdash e_{\text{arr}}[e_{\text{idx}}] \Rightarrow T} \\[0.8em]
\text{\bf (T-Struct-Field)} \quad & \frac{\Gamma \vdash e_{\text{str}} \Rightarrow S \quad \text{field\_type}(S, f) = \tau}{\Gamma \vdash e_{\text{str}}.f \Rightarrow \tau} \\[0.8em]
\text{\bf (T-Enum-Cons)} \quad & \frac{\text{variant}(E, K) = (\tau_1, \dots, \tau_m) \quad \forall i.\; \Gamma \vdash e_i \Leftarrow \tau_i}{\Gamma \vdash E::K(e_1, \dots, e_m) \Rightarrow E}
\end{align}

\section{The Maranget Pattern Matching Algorithm}
\label{sec:types:maranget}

NumLang validates pattern match exhaustiveness using Luc Maranget's matrix-based algorithm \citep{maranget2007warnings}. Patterns in a match expression are structured into an $n \times m$ pattern matrix $P$, where each row represents a match arm pattern tuple and each column corresponds to an inspected value.

Algorithm~\ref{alg:maranget} defines the utility predicate $\mathcal{U}(P, \vec{q})$ which determines whether pattern vector $\vec{q}$ is useful relative to matrix $P$.

\begin{algorithm}[h!]
\caption{Maranget Pattern Matrix Utility Algorithm}
\label{alg:maranget}
\begin{algorithmic}[1]
\Require Pattern Matrix $P$, Target Pattern Vector $\vec{q} = (q_1, \dots, q_k)$
\Ensure Boolean indicating whether $\vec{q}$ matches values not covered by $P$
\Procedure{IsUseful}{$P, \vec{q}$}
    \If{$P$ has $0$ rows} \State \Return \textbf{true} \EndIf
    \If{$k == 0$} \State \Return \textbf{false} \EndIf
    \If{$q_1$ is a constructor $C(r_1, \dots, r_a)$}
        \State Let $S(C, P)$ be the specialized matrix projecting constructor $C$
        \State \Return \Call{IsUseful}{$S(C, P), (r_1, \dots, r_a, q_2, \dots, q_k)$}
    \ElsIf{$q_1$ is a wildcard $\_$}
        \If{The set of constructors appearing in column 1 of $P$ is complete}
            \ForAll{constructors $C \in \operatorname{Constructors}(\text{type}(q_1))$}
                \If{\Call{IsUseful}{$S(C, P), (\dots, q_2, \dots, q_k)$}}
                    \State \Return \textbf{true}
                \EndIf
            \EndFor
            \State \Return \textbf{false}
        \Else
            \State Let $D(P)$ be the default matrix containing rows with wildcards
            \State \Return \Call{IsUseful}{$D(P), (q_2, \dots, q_k)$}
        \EndIf
    \EndIf
\EndProcedure
\end{algorithmic}
\end{algorithm}

\section{Bidirectional Type Inference Engine Implementation}
\label{sec:types:checker_impl}

The type inference engine implemented in \texttt{src/typecheck/checker.rs} evaluates expressions through bidirectional typing. Listing~\ref{lst:checker_impl} presents the core expression inference routines, environment management, and unification logic:

\begin{lstlisting}[language=Rust, caption={Core Bidirectional Type Checking and Unification (\texttt{src/typecheck/checker.rs}).}, label={lst:checker_impl}]
        span: Span,
    },

    #[error("Condition must evaluate to bool, found {found}")]
    InvalidConditionType {
        found: Type,
        span: Span,
    },

    #[error("Invalid binary operands for '{op:?}': left is {left}, right is {right}")]
    InvalidBinaryOperands {
        op: BinaryOp,
        left: Type,
        right: Type,
        span: Span,
    },

    #[error("Invalid unary operand for '{op:?}': found {found}")]
    InvalidUnaryOperand {
        op: UnaryOp,
        found: Type,
        span: Span,
    },

    #[error("Function '{name}' expected {expected} arguments, but received {found}")]
    ArityMismatch {
        name: String,
        expected: usize,
        found: usize,
        span: Span,
    },

    #[error("Unknown type '{name}'")]
    UnknownType {
        name: String,
        span: Span,
    },

    #[error("Function return type mismatch: expected {expected}, found {found}")]
    InvalidReturn {
        expected: Type,
        found: Type,
        span: Span,
    },

    #[error("Cannot index non-array type '{found}'")]
    CannotIndexNonArray {
        found: Type,
        span: Span,
    },

    #[error("Array index must be an integer, found '{found}'")]
    InvalidIndexType {
        found: Type,
        span: Span,
    },

    #[error("Array literal cannot be empty")]
    EmptyArrayLiteral {
        span: Span,
    },

    #[error("Array index {index} out of bounds for array of length {len}")]
    IndexOutOfBounds {
        index: i64,
        len: usize,
        span: Span,
    },

    #[error("Array element type mismatch: expected {expected}, found {found}")]
    ArrayElementMismatch {
        expected: Type,
        found: Type,
        span: Span,
    },

    #[error("'break' may only be used inside a loop")]
    BreakOutsideLoop { span: Span },

    #[error("'continue' may only be used inside a loop")]
    ContinueOutsideLoop { span: Span },

    #[error("Struct '{name}' has no field '{field}'")]
    NoSuchField {
        name: String,
        field: String,
        span: Span,
    },

    #[error("Missing field '{field}' in struct '{name}' initialization")]
    MissingField {
        name: String,
        field: String,
        span: Span,
    },

    #[error("Cannot access field on non-struct type '{found}'")]
    CannotAccessFieldNonStruct {
        found: Type,
        span: Span,
    },

    #[error("Non-exhaustive match: missing wildcard '_' or uncovered cases")]
    NonExhaustiveMatch { span: Span },

    #[error("Match expression must have at least one arm")]
    EmptyMatch { span: Span },

    #[error("Enum '{enum_name}' has no variant '{variant_name}'")]
    NoSuchVariant {
        enum_name: String,
        variant_name: String,
        span: Span,
    },

    #[error("Cannot match variant pattern on non-enum type '{found}'")]
    CannotMatchNonEnum {
        found: Type,
        span: Span,
    },

    #[error("Variant '{enum_name}::{variant_name}' expected {expected} payload arguments, found {found}")]
    PayloadArityMismatch {
        enum_name: String,
        variant_name: String,
        expected: usize,
        found: usize,
        span: Span,
    },

    #[error("Type '{ty}' is not callable")]
    NotCallable {
        ty: Type,
        span: Span,
    },
}

impl TypeError {
    pub fn span(&self) -> Span {
        match self {
            TypeError::TypeMismatch { span, .. } => *span,
            TypeError::UndeclaredVariable { span, .. } => *span,
            TypeError::UndeclaredFunction { span, .. } => *span,
            TypeError::CannotMutateImmutable { span, .. } => *span,
            TypeError::DuplicateDeclaration { span, .. } => *span,
            TypeError::InvalidConditionType { span, .. } => *span,
            TypeError::InvalidBinaryOperands { span, .. } => *span,
            TypeError::InvalidUnaryOperand { span, .. } => *span,
            TypeError::ArityMismatch { span, .. } => *span,
            TypeError::UnknownType { span, .. } => *span,
            TypeError::InvalidReturn { span, .. } => *span,
            TypeError::CannotIndexNonArray { span, .. } => *span,
            TypeError::InvalidIndexType { span, .. } => *span,
            TypeError::EmptyArrayLiteral { span, .. } => *span,
            TypeError::IndexOutOfBounds { span, .. } => *span,
            TypeError::ArrayElementMismatch { span, .. } => *span,
            TypeError::BreakOutsideLoop { span } => *span,
            TypeError::ContinueOutsideLoop { span } => *span,
            TypeError::NoSuchField { span, .. } => *span,
            TypeError::MissingField { span, .. } => *span,
            TypeError::CannotAccessFieldNonStruct { span, .. } => *span,
            TypeError::NonExhaustiveMatch { span } => *span,
            TypeError::EmptyMatch { span } => *span,
            TypeError::NoSuchVariant { span, .. } => *span,
            TypeError::CannotMatchNonEnum { span, .. } => *span,
            TypeError::PayloadArityMismatch { span, .. } => *span,
            TypeError::NotCallable { span, .. } => *span,
        }
    }

    pub fn error_code(&self) -> &'static str {
        match self {
            TypeError::TypeMismatch { .. } => "E001",
            TypeError::UndeclaredVariable { .. } => "E002",
            TypeError::UndeclaredFunction { .. } => "E003",
            TypeError::CannotMutateImmutable { .. } => "E004",
            TypeError::DuplicateDeclaration { .. } => "E005",
            TypeError::InvalidConditionType { .. } => "E006",
            TypeError::InvalidBinaryOperands { .. } => "E007",
            TypeError::InvalidUnaryOperand { .. } => "E008",
            TypeError::ArityMismatch { .. } => "E009",
            TypeError::UnknownType { .. } => "E010",
            TypeError::InvalidReturn { .. } => "E011",
            TypeError::CannotIndexNonArray { .. } => "E012",
            TypeError::InvalidIndexType { .. } => "E013",
            TypeError::EmptyArrayLiteral { .. } => "E014",
            TypeError::IndexOutOfBounds { .. } => "E015",
            TypeError::ArrayElementMismatch { .. } => "E016",
            TypeError::BreakOutsideLoop { .. } => "E017",
            TypeError::ContinueOutsideLoop { .. } => "E018",
            TypeError::NoSuchField { .. }
            | TypeError::MissingField { .. }
            | TypeError::CannotAccessFieldNonStruct { .. } => "E019",
            TypeError::NonExhaustiveMatch { .. }
            | TypeError::EmptyMatch { .. } => "E020",
            TypeError::NoSuchVariant { .. }
            | TypeError::CannotMatchNonEnum { .. }
            | TypeError::PayloadArityMismatch { .. } => "E021",
            TypeError::NotCallable { .. } => "E022",
        }
    }
}

use std::collections::HashMap;

#[derive(Debug, Clone)]
pub struct StructInfo {
    pub name: String,
    pub fields: Vec<(String, Type)>,
    pub field_indices: HashMap<String, usize>,
    pub field_offsets: HashMap<String, u32>,
    pub total_size: u32,
    pub align: u32,
    pub span: Span,
}

#[derive(Debug, Clone)]
pub struct EnumInfo {
    pub name: String,
    pub variants: Vec<EnumVariantInfo>,
    pub variant_indices: HashMap<String, usize>,
    pub max_payload_size: u32,
    pub span: Span,
}

#[derive(Debug, Clone)]
pub struct EnumVariantInfo {
    pub name: String,
    pub tag: usize,
    pub payload: Vec<Type>,
    pub span: Span,
}

pub struct TypeChecker {
    env: ScopeEnvironment,
    current_fn_return_ty: Type,
    active_loop_bounds: HashMap<String, i64>,
    loop_depth: usize,
    pub struct_infos: HashMap<String, StructInfo>,
    pub enum_infos: HashMap<String, EnumInfo>,
    pub variant_to_enum: HashMap<String, Vec<String>>,
    pub current_type_params: Vec<String>,
}

impl Default for TypeChecker {
    fn default() -> Self {
        Self::new()
    }
}


\end{lstlisting}

\section{Whole-Program Monomorphization Implementation}
\label{sec:types:monomorphize_impl}

NumLang achieves zero-cost generics through static monomorphization. The monomorphization pass in \texttt{src/opt/monomorphize.rs} traces generic call graphs, computes concrete type instantiation maps, and synthesizes specialized function clones prior to intermediate representation lowering:

\begin{lstlisting}[language=Rust, caption={Whole-Program Monomorphization Engine (\texttt{src/opt/monomorphize.rs}).}, label={lst:monomorphize_impl}]
//!
//! After type-checking, any function with `type_params` is a template.
//! This pass scans all call sites, collects the concrete type arguments,
//! and produces a concrete copy of each generic function for each
//! distinct instantiation. Generic originals are removed from the program.

use std::collections::HashMap;
use crate::span::Span;
use crate::typecheck::checker::TypeError;
use crate::typecheck::typed_ast::{
    TypedBlock, TypedExpr, TypedFunction, TypedLiteral, TypedMatchPattern, TypedProgram,
    TypedStmt,
};
use crate::typecheck::types::{ClosureType, Type};

/// Substitution: map from type param name to concrete Type
pub type Subst = HashMap<String, Type>;

pub fn substitute_type(ty: &Type, subst: &Subst) -> Type {
    match ty {
        Type::Param(name) => subst.get(name).cloned().unwrap_or_else(|| ty.clone()),
        Type::Array(elem, len) => Type::Array(Box::new(substitute_type(elem, subst)), *len),
        Type::Box(inner) => Type::Box(Box::new(substitute_type(inner, subst))),
        Type::Ptr(inner) => Type::Ptr(Box::new(substitute_type(inner, subst))),
        Type::Fn(args, ret) => Type::Fn(
            args.iter().map(|a| substitute_type(a, subst)).collect(),
            Box::new(substitute_type(ret, subst)),
        ),
        Type::Closure(c) => Type::Closure(Box::new(ClosureType {
            params: c.params.iter().map(|p| substitute_type(p, subst)).collect(),
            ret: Box::new(substitute_type(&c.ret, subst)),
            captured: c
                .captured
                .iter()
                .map(|(n, t)| (n.clone(), substitute_type(t, subst)))
                .collect(),
        })),
        other => other.clone(),
    }
}

pub fn unify_types(
    param_ty: &Type,
    arg_ty: &Type,
    subst: &mut Subst,
    span: Span,
) -> Result<(), TypeError> {
    match (param_ty, arg_ty) {
        (Type::Param(name), actual) => {
            if let Some(existing) = subst.get(name) {
                if existing != actual && !actual.is_compatible_with(existing) {
                    return Err(TypeError::TypeMismatch {
                        expected: existing.clone(),
                        found: actual.clone(),
                        span,
                    });
                }
            } else {
                subst.insert(name.clone(), actual.clone());
            }
            Ok(())
        }
        (Type::Array(p_elem, p_len), Type::Array(a_elem, a_len)) if p_len == a_len => {
            unify_types(p_elem, a_elem, subst, span)
        }
        (Type::Box(p_inner), Type::Box(a_inner)) | (Type::Ptr(p_inner), Type::Ptr(a_inner)) => {
            unify_types(p_inner, a_inner, subst, span)
        }
        (Type::Fn(p_args, p_ret), Type::Fn(a_args, a_ret)) => {
            if p_args.len() != a_args.len() {
                return Err(TypeError::TypeMismatch {
                    expected: param_ty.clone(),
                    found: arg_ty.clone(),
                    span,
                });
            }
            for (pa, aa) in p_args.iter().zip(a_args) {
                unify_types(pa, aa, subst, span)?;
            }
            unify_types(p_ret, a_ret, subst, span)
        }
        (Type::Fn(p_args, p_ret), Type::Closure(c)) => {
            if p_args.len() != c.params.len() {
                return Err(TypeError::TypeMismatch {
                    expected: param_ty.clone(),
                    found: arg_ty.clone(),
                    span,
                });
            }
            for (pa, ca) in p_args.iter().zip(&c.params) {
                unify_types(pa, ca, subst, span)?;
            }
            unify_types(p_ret, &c.ret, subst, span)
        }
        (expected, found) => {
            if expected.is_compatible_with(found) {
                Ok(())
            } else {
                Err(TypeError::TypeMismatch {
                    expected: expected.clone(),
                    found: found.clone(),
                    span,
                })
            }
        }
    }
}

pub fn type_to_ident(ty: &Type) -> String {
    match ty {
        Type::I8 => "i8".to_string(),
        Type::I16 => "i16".to_string(),
        Type::I32 => "i32".to_string(),
        Type::I64 => "i64".to_string(),
        Type::U8 => "u8".to_string(),
        Type::U16 => "u16".to_string(),
        Type::U32 => "u32".to_string(),
        Type::U64 => "u64".to_string(),
        Type::Usize => "usize".to_string(),
        Type::F32 => "f32".to_string(),
        Type::F64 => "f64".to_string(),
        Type::Bool => "bool".to_string(),
        Type::Void => "void".to_string(),
        Type::Str => "str".to_string(),
        Type::Struct(s) => s.clone(),
        Type::Enum(e) => e.clone(),
        Type::Array(elem, len) => format!("arr_{}_{}", type_to_ident(elem), len),
        Type::Box(inner) => format!("box_{}", type_to_ident(inner)),
        Type::Ptr(inner) => format!("ptr_{}", type_to_ident(inner)),
        Type::Fn(args, ret) => format!(
            "fn_{}_ret_{}",
            args.iter().map(type_to_ident).collect::<Vec<_>>().join("_"),
            type_to_ident(ret)
        ),
        Type::Closure(c) => format!(
            "closure_{}_ret_{}",
            c.params.iter().map(type_to_ident).collect::<Vec<_>>().join("_"),
            type_to_ident(&c.ret)
        ),
        Type::Param(p) => p.clone(),
    }
}

pub fn mangle_generic_name(base: &str, subst: &Subst, type_params: &[String]) -> String {
    let mut parts = Vec::new();
    for tp in type_params {
        if let Some(ty) = subst.get(tp) {
            parts.push(format!("{}_{}", tp, type_to_ident(ty)));
        }
    }
    format!("{}__{}", base, parts.join("__"))
}

pub fn substitute_expr(expr: &mut TypedExpr, subst: &Subst) {
    match expr {
        TypedExpr::Literal { lit, ty, .. } => {
            *ty = substitute_type(ty, subst);
            if let TypedLiteral::Int(n, lty) = lit {
                *lty = substitute_type(lty, subst);
                let _ = n;
            }
        }
        TypedExpr::Ident { ty, .. } => {
            *ty = substitute_type(ty, subst);
        }
        TypedExpr::Unary { expr, ty, .. } => {
            *ty = substitute_type(ty, subst);
            substitute_expr(expr, subst);
        }
        TypedExpr::Binary { left, right, ty, .. } => {
            *ty = substitute_type(ty, subst);
            substitute_expr(left, subst);
            substitute_expr(right, subst);
        }
        TypedExpr::Call { args, ty, .. } => {
            *ty = substitute_type(ty, subst);
            for arg in args {
                substitute_expr(arg, subst);
            }
        }
        TypedExpr::ArrayLiteral { elements, ty, .. } => {
            *ty = substitute_type(ty, subst);
            for el in elements {
                substitute_expr(el, subst);
            }
        }
        TypedExpr::Index { target, index, ty, .. } => {
            *ty = substitute_type(ty, subst);
            substitute_expr(target, subst);
            substitute_expr(index, subst);
        }
        TypedExpr::StructLiteral { fields, ty, .. } => {
            *ty = substitute_type(ty, subst);
            for (_, f_expr) in fields {
                substitute_expr(f_expr, subst);
            }
        }
        TypedExpr::FieldAccess { target, ty, .. } => {
            *ty = substitute_type(ty, subst);
            substitute_expr(target, subst);

\end{lstlisting}
